\begin{mistake}{2026-04-20}{03}

\begin{problemcontent}
设 \(\mathbf{A}, \mathbf{B}\) 为 \(n\) 阶方阵，\(\mathbf{P}, \mathbf{Q}\) 为 \(n\) 阶可逆矩阵，下列命题不正确的是（\quad）

(A) 若 \(\mathbf{B} = \mathbf{A}\mathbf{Q}\)，则 \(\mathbf{A}\) 的列向量组与 \(\mathbf{B}\) 的列向量组等价.

(B) 若 \(\mathbf{B} = \mathbf{P}\mathbf{A}\)，则 \(\mathbf{A}\) 的行向量组与 \(\mathbf{B}\) 的行向量组等价.

(C) 若 \(\mathbf{B} = \mathbf{P}\mathbf{A}\mathbf{Q}\)，则 \(\mathbf{A}\) 的行(列)向量组与 \(\mathbf{B}\) 的行(列)向量组等价.

(D) 若 \(\mathbf{A}\) 的行(列)向量组与矩阵 \(\mathbf{B}\) 的行(列)向量组等价，则矩阵 \(\mathbf{A}\) 与 \(\mathbf{B}\) 等价.
\end{problemcontent}

\begin{solutioncontent}
本题选 (C)。

(A) \(\mathbf{B} = \mathbf{A}\mathbf{Q} \implies \mathbf{B}\) 的每一列均为 \(\mathbf{A}\) 的列向量的线性组合。因为 \(\mathbf{Q}\) 可逆，故 \(\mathbf{A} = \mathbf{B}\mathbf{Q}^{-1}\)，同理 \(\mathbf{A}\) 的每一列也可由 \(\mathbf{B}\) 的列向量线性表示。两者互相线性表示，列向量组等价，命题正确。

(B) 同理，\(\mathbf{B} = \mathbf{P}\mathbf{A}\) 且 \(\mathbf{P}\) 可逆，\(\mathbf{A}\) 与 \(\mathbf{B}\) 的行向量组可互相线性表示，等价，命题正确。

(D) 向量组等价 \(\implies\) 向量组的秩相等 \(\implies\) 矩阵的秩相等。对于同型方阵，秩相等即为矩阵等价的充要条件，命题正确。

(C) \(\mathbf{B} = \mathbf{P}\mathbf{A}\mathbf{Q}\) 仅表示矩阵等价（秩相等），不能保证生成空间相同。
反例：设 \(\mathbf{A}=\begin{pmatrix} 1 & 0 \\ 0 & 0 \end{pmatrix}, \mathbf{P}=\begin{pmatrix} 0 & 1 \\ 1 & 0 \end{pmatrix}, \mathbf{Q}=\begin{pmatrix} 1 & 0 \\ 0 & 1 \end{pmatrix}\)，则 \(\mathbf{B}=\begin{pmatrix} 0 & 0 \\ 1 & 0 \end{pmatrix}\)。\(\mathbf{A}\) 的列空间在 \(x\) 轴，\(\mathbf{B}\) 的列空间在 \(y\) 轴，两者不能互相表示，列向量组不等价。命题不正确。
\end{solutioncontent}

\mistakeknowledge{
\begin{itemize}
 \item \textbf{核心公式/定理}：同型方阵 \(\mathbf{A}\) 与 \(\mathbf{B}\) 等价 \(\iff r(\mathbf{A}) = r(\mathbf{B})\)。向量组等价必然导致二者秩相等。
 \item \textbf{易错点}：混淆“矩阵等价”与“向量组等价”。矩阵等价是弱条件（仅要求秩等），向量组等价是强条件（要求底层张成的线性空间完全重合）。
 \item \textbf{关联知识}：“左乘行，右乘列”。对矩阵左乘可逆阵，行空间不变（行向量组等价）；右乘可逆阵，列空间不变（列向量组等价）。
\end{itemize}}

\end{mistake}
